\documentclass[12pt]{article}

\usepackage[T1]{fontenc}
\usepackage[utf8]{inputenc}
\usepackage[english]{babel}

\usepackage{lmodern}
\usepackage{geometry}
\usepackage{graphicx}
\usepackage{fancyhdr}
\usepackage{amsmath}
\usepackage{tocloft}
\usepackage{titlesec}
\usepackage{hyperref}
\usepackage{float}
\usepackage{listings}
\usepackage{xcolor}
\usepackage{tocloft}
% --- Fix TOC number box widths for Roman numerals ---
\setlength{\cftsecnumwidth}{3.0em}     % Gives enough space for 'VII'
\setlength{\cftsubsecnumwidth}{4.0em}  % Gives enough space for 'VII.ii'
\setlength{\cftsubsubsecnumwidth}{4.8em}
\hypersetup{
    colorlinks=true,
    linkcolor=black,
    urlcolor=black
}

\renewcommand\thesection{\Roman{section}}
\renewcommand\thesubsection{\thesection.\roman{subsection}}

\geometry{
    a4paper, top=1in, bottom=1in, left=1in, right=1in,
    headheight=15pt, footskip=15pt,
}

\lstset{
    language=Matlab, basicstyle=\ttfamily\small,
    keywordstyle=\color{blue}, commentstyle=\color{gray}, stringstyle=\color{red},
    numbers=left, numberstyle=\tiny, stepnumber=1, numbersep=5pt,
    backgroundcolor=\color{white}, showspaces=false, showstringspaces=false,
    showtabs=false, frame=single, tabsize=4, captionpos=b, breaklines=true
}

\linespread{1.06}
\setlength{\parindent}{0pt}
\setlength{\parskip}{0.8em}

\pagestyle{fancy}
\fancyhf{}
\lhead{DVD-ML-Driven PVT Variation Aware Leakage Estimation}
\cfoot{\thepage}
\renewcommand{\headrulewidth}{0.4pt}
\renewcommand{\footrulewidth}{0pt}

\titlespacing*{\section}{0pt}{8pt}{6pt}
\titlespacing*{\subsection}{0pt}{6pt}{4pt}

\newcommand{\includelogo}{
    \IfFileExists{IIIT_Hyderabad_Logo.jpg}{%
        \includegraphics[width=.8\textwidth]{IIIT_Hyderabad_Logo.jpg}%
    }{%
        \fbox{\parbox[c][4cm][c]{0.8\textwidth}{\centering\Large Image placeholder\\[0.5em]\small IIIT Hyderabad Logo}}%
    }%
}

\newcommand{\includeimage}[3][0.75\linewidth]{%
    \IfFileExists{#2}{%
        \includegraphics[width=#1]{#2}%
    }{%
        \fbox{%
            \parbox[c][4.5cm][c]{#1}{%
                \centering
                \Large Image Placeholder\\[0.5em]
                \small #3 (\texttt{#2})
            }%
        }%
    }%
}

% ---------------------------------------------------------------------------
% Numbers and tables are generated from the project outputs by
% make_report_assets.py (run: python report/make_report_assets.py).
% ---------------------------------------------------------------------------
% generated by make_report_assets.py - do not edit
\newcommand{\AugAfterMin}{98.27}
\newcommand{\AugOnAfterMax}{99.05}
\newcommand{\AugOnAfterMin}{98.57}
\newcommand{\AugOnBeforeMax}{88.57}
\newcommand{\AugOnBeforeMin}{66.96}
\newcommand{\ClaBias}{-0.318}
\newcommand{\ClaCompMeanAcc}{99.64}
\newcommand{\ClaCompMedErr}{0.241}
\newcommand{\ClaErrCold}{0.290}
\newcommand{\ClaErrHot}{1.220}
\newcommand{\ClaMaxErr}{4.74}
\newcommand{\ClaMeanAcc}{99.31}
\newcommand{\ClaMedErr}{0.500}
\newcommand{\ClaModelMeanAcc}{99.51}
\newcommand{\ClaModelMedErr}{0.357}
\newcommand{\ClaPninetyfive}{1.91}
\newcommand{\GateBest}{NOR2 11 (median 0.097\,\%)}
\newcommand{\GateBestGate}{NOR2}
\newcommand{\GateBias}{+0.030}
\newcommand{\GateErrCold}{0.346}
\newcommand{\GateErrHot}{0.469}
\newcommand{\GateMaxErr}{6.12}
\newcommand{\GateMeanAcc}{99.42}
\newcommand{\GateMedErr}{0.397}
\newcommand{\GateNinetyNineCount}{17}
\newcommand{\GatePninetyfive}{1.76}
\newcommand{\GateWorst}{OR2 00}
\newcommand{\GateWorstAcc}{98.59}
\newcommand{\GateWorstBias}{-1.18}
\newcommand{\GateWorstGate}{INV}
\newcommand{\GateWorstPninetyfive}{3.42}
\newcommand{\KclWorstAbs}{$9.85\times10^{-13}$}
\newcommand{\KclWorstRel}{$1.22\times10^{-4}$}
\newcommand{\ModelMaxAcc}{99.95}
\newcommand{\ModelMinAcc}{98.93}
\newcommand{\ModelMinAccWhere}{NMOS ON ID}
\newcommand{\NmosOnShare}{2}
\newcommand{\OffRowsN}{614}
\newcommand{\OffRowsP}{635}
\newcommand{\OnShare}{42}
\newcommand{\OnTrainSmallVds}{2.2}
\newcommand{\OnVdsMedian}{0.49}
\newcommand{\OnVdsNinety}{3.1}
\newcommand{\PmosOffShare}{58}
\newcommand{\PmosOnShare}{39}
\newcommand{\RhoTempN}{0.82}
\newcommand{\RhoTempP}{0.72}
\newcommand{\RhoToxMax}{-0.43}
\newcommand{\RhoToxMin}{-0.60}
\newcommand{\RhoVthMax}{-0.33}
\newcommand{\RhoVthMin}{-0.54}

\newcommand{\gentable}[1]{\resizebox{\linewidth}{!}{\input{generated/#1.tex}}}
\newcommand{\gentablenat}[1]{\input{generated/#1.tex}}

\begin{document}

\begin{titlepage}
    \centering
    \vspace*{12mm}
    \includelogo
    \vspace{10mm}
    {\large\scshape International Institute of Information Technology\\ Hyderabad \par}
    \vspace{6mm}
    {\large \bfseries Digital VLSI Design \par}
    \vspace{8pt}
    {\large \bfseries ML-Driven PVT Variation Aware Leakage Estimation  \par}
    \vspace{6pt}
    {\large \textit{Course Code: M26EC2.408} \par}
    \vfill
    {\Large\scshape \bfseries Anumay Rai \& Ravi Kumar Sahu	 \& Diddigi Srivatsa
    \& Ashlin K Sunil \par}
    \vspace{4pt}
    {\large Roll No: 2024102024 \& 2025122009 \par}
    \vspace{10mm}
    \thispagestyle{empty}
\end{titlepage}

\clearpage
\pagenumbering{roman}
\setcounter{page}{1}
\thispagestyle{plain}
\tableofcontents
\clearpage

\pagenumbering{arabic}
\setcounter{page}{1}
\pagestyle{fancy}
\fontsize{11}{13}\selectfont

% ===========================================================================
\section{Introduction}
% ===========================================================================
Leakage power in deep sub-micron CMOS depends strongly on process, supply voltage and temperature
(PVT). Estimating it with SPICE at every corner and input vector is accurate but slow. In this
project we replace SPICE by machine-learning (ML) models of a \emph{single transistor} and build
gate and circuit leakage hierarchically from them:

\begin{enumerate}
    \item \textbf{Part A -- transistor models.} SPICE data for NMOS and PMOS in their ON and OFF
          states, and four models that predict all four terminal currents $(I_d, I_g, I_s, I_b)$.
    \item \textbf{Part B -- gates.} INV, NAND2, NOR2, AND2 and OR2 for every input combination at
          200 random corners; leakage predicted only from the four transistor models and compared
          with SPICE of the gate netlist.
    \item \textbf{Part C -- test circuit.} A 3-bit carry-lookahead adder (CLA) composed from the gate
          predictions and compared with SPICE of the complete netlist.
\end{enumerate}

Throughout, the error of one prediction is $e = |I_{ML}-I_{SPICE}|/|I_{SPICE}|$ (in \%), and the
\emph{mean accuracy} is $100 - \operatorname{mean}(\min(e, 100\,\%))$. The target is a mean
accuracy of up to 99\,\%.

% ===========================================================================
\section{Setup}
% ===========================================================================
\begin{table}[H]
\centering
\begin{tabular}{ll}
\hline
Item & Value \\
\hline
Model file & PTM \texttt{22nm\_MGK.pm} (BSIM4, metal gate / high-$k$) \\
Nominal supply & $V_{DD} = 0.9$\,V, sampled $0.81$--$0.99$\,V ($\pm10\,\%$) \\
Channel length & $L = L_{min} = 22$\,nm \\
Widths (Part A) & $W = 1$--$8\times L_{min}$ = 22--176\,nm \\
Widths (Parts B, C) & $W_n = 2L_{min} = 44$\,nm, $W_p = 4L_{min} = 88$\,nm \\
Temperature & $-40$ to $125\,^\circ$C \\
Process parameters & \texttt{vth0, toxe, ndep, u0, xj}: nominal $\pm10\,\%$, separate for NMOS and PMOS \\
Test circuit & 3-bit carry-lookahead adder (39 gates, 186 transistors) \\
Simulator & ngspice-42, DC operating point \\
SPICE options & \texttt{gmin=1e-18 abstol=1e-18 reltol=1e-6 vntol=1e-9} \\
ML software & Python 3.10, LightGBM, scikit-learn, SciPy \\
\hline
\end{tabular}
\caption{Assignment and tools.}
\end{table}

The nominal process values were read from the model card: NMOS \texttt{vth0}\,=\,0.3692\,V,
\texttt{toxe}\,=\,0.65\,nm, \texttt{ndep}\,=\,$1.2\times10^{19}$, \texttt{u0}\,=\,0.0181,
\texttt{xj}\,=\,7.2\,nm; PMOS \texttt{vth0}\,=\,$-0.25399$\,V, \texttt{toxe}\,=\,0.67\,nm,
\texttt{ndep}\,=\,$4.4\times10^{18}$, \texttt{u0}\,=\,0.0023, \texttt{xj}\,=\,7.2\,nm.

% ===========================================================================
\section{Data Generation}
% ===========================================================================
\subsection{Sampling}
For each of the four cases 10{,}000 samples were drawn with a fixed seed (42). $V_{DD}$ is drawn
first, then $V_d$ and $V_s$ independently and uniformly in $[0, V_{DD}]$, the width, the temperature
and the five process parameters, all uniform and independent. The gate and body voltages are fixed
by the case:

\begin{table}[H]
\centering
\begin{tabular}{lcc}
\hline
Case & $V_g$ & $V_b$ \\
\hline
NMOS ON  & $V_{DD}$ & 0 \\
NMOS OFF & 0 & 0 \\
PMOS ON  & 0 & $V_{DD}$ \\
PMOS OFF & $V_{DD}$ & $V_{DD}$ \\
\hline
\end{tabular}
\caption{Gate and body bias per case.}
\end{table}

\subsection{Netlists and simulation}
Each sample is a single-transistor netlist in which every terminal has its own DC voltage source
(\texttt{Vd, Vg, Vs, Vb}) and the required \texttt{.options} line is set. To keep run time low, 1{,}000
samples are simulated per ngspice run: the netlist is loaded once and, for each sample, the process
parameters are changed with \texttt{altermod}, the width and voltages with \texttt{alter} and the
temperature with \texttt{option temp}, followed by an \texttt{op} analysis. The terminal currents are
the currents of the four sources with the sign flipped, so a positive current flows \emph{into} the
transistor terminal. ngspice also writes back every input it actually simulated; these read-back
values are what the dataset stores, so the inputs and outputs always belong together
(read-back mismatch $<10^{-8}$). One example netlist per case
(\texttt{dataset/<CASE>/example\_<case>.net}) reproduces its dataset row to $<0.001\,\%$.

To validate the dataset, 300 random rows of every case were re-simulated independently; the stored
currents were reproduced to a median of $\sim10^{-6}\,\%$.

\subsection{Train / validation / test split}
Each dataset is split randomly 70\,/\,15\,/\,15\,\% (7{,}000 / 1{,}500 / 1{,}500 samples) with the
fixed seed 42. The same split is used by every model. Physics formulas and ML models are fitted on
the training set only; early stopping, blend weights and model selection use the validation set;
the test set is used once, for the final numbers.

\subsection{Data check}
\textbf{Kirchhoff check.} Table~\ref{tab:kcl} lists $|I_d+I_g+I_s+I_b|$ for every case. All
40{,}000 rows pass (relative residual $<10^{-6}$ or absolute residual $<10^{-14}$\,A); the worst
relative residual is \KclWorstRel{} and the worst absolute residual \KclWorstAbs\,A.

\begin{table}[H]
\centering
\gentable{tab_kcl}
\caption{Kirchhoff check per case (relative = residual / largest terminal current).}
\label{tab:kcl}
\end{table}

\textbf{Trend plots.} The OFF current is the subthreshold channel current per width, measured at the
source-side terminal (where gate tunnelling is negligible) at the textbook OFF bias: source-side
terminal within 0.1\,V of its rail ($V_{gs}\approx0$) and $|V_{ds}|\ge0.6$\,V
(\OffRowsN{} NMOS and \OffRowsP{} PMOS samples).

\begin{itemize}
    \item \textbf{OFF current vs temperature} (Fig.~\ref{fig:trend_temp}): rises by one to two decades
          from $-40$ to $125\,^\circ$C, as expected from the $\exp(-V_{th}/nV_T)$ dependence and the
          drop of $V_{th}$ with temperature (Spearman $\rho$ = \RhoTempN{} NMOS, \RhoTempP{} PMOS).
    \item \textbf{OFF current vs $|$vth0$|$} (Fig.~\ref{fig:trend_vth}): at a fixed temperature band the
          OFF current falls with $|$vth0$|$ ($\rho$ between \RhoVthMin{} and \RhoVthMax), the
          subthreshold exponential in $V_{th}$. Temperature dominates the overall spread, so the
          plot shows one line per temperature band.
    \item \textbf{Gate current vs toxe} (Fig.~\ref{fig:trend_tox}): $|I_g|/W$ falls exponentially
          with oxide thickness in all four cases ($\rho$ between \RhoToxMin{} and \RhoToxMax), the
          signature of direct tunnelling. ON devices have the larger gate current because the whole
          channel is inverted.
\end{itemize}

\begin{figure}[H]
\centering
\includeimage[0.72\linewidth]{figures/trend_off_current_vs_temp.png}{OFF current vs temperature}
\caption{OFF current vs temperature (NMOS OFF, PMOS OFF). Dots: samples; lines: median per bin.}
\label{fig:trend_temp}
\end{figure}

\begin{figure}[H]
\centering
\includeimage[0.95\linewidth]{figures/trend_off_current_vs_vth0.png}{OFF current vs vth0}
\caption{OFF current vs $|$vth0$|$, one median line per temperature band.}
\label{fig:trend_vth}
\end{figure}

\begin{figure}[H]
\centering
\includeimage[0.72\linewidth]{figures/trend_gate_current_vs_toxe.png}{Gate current vs toxe}
\caption{Gate current per width vs oxide thickness, all four cases.}
\label{fig:trend_tox}
\end{figure}

% ===========================================================================
\section{Transistor Models}
% ===========================================================================
\subsection{Approach: physics baseline + LightGBM correction}
All four chosen models share one structure. A compact physics model is fitted first (SciPy
\texttt{least\_squares}, training rows only) and LightGBM learns only the small correction on top:

\begin{itemize}
    \item \textbf{Physics baseline.} Subthreshold / EKV-style channel current $C$ (body effect, DIBL,
          temperature, width and process terms, plus a smooth quadratic correction), gate tunnelling
          at the drain and source edges $F_d, F_s$, and a body-current model (junction, high-temperature
          and gate-to-body terms). The physics guess is $I_d=\mathrm{sgn}(V_d-V_s)\,C+F_d$,
          $I_s=-\mathrm{sgn}(V_d-V_s)\,C+F_s$, $I_g=-(F_d+F_s)$.
    \item \textbf{Targets (preprocessing).} For $I_d$ and $I_s$ the model learns
          $I/(C+F_\text{this})$ minus its physics value. This ratio is smooth through $I=0$ and carries
          the sign, which removes the large relative errors a $\log|I|$ target produces where the
          current changes sign. For $I_g$ and $I_b$ the target is $\ln|I| - \ln|I_\text{phys}|$.
    \item \textbf{One model for all four currents.} The four currents of a sample are stacked as four
          rows with a categorical ``current'' feature; $I_d$ and $I_s$ rows are written from the
          view of ``this terminal'' and ``the other terminal'', so one function learns both.
    \item \textbf{Features.} The raw inputs plus $\log W$, $V_{high}$, $V_{low}$, $|V_{ds}|$,
          $1/V_T$, $V_{th}/V_T$, $|V_{ds}|/V_T$ and the physics quantities ($\log C$, $\log F$, share
          of the physics current).
    \item \textbf{Mirroring.} NMOS ON is the mirror image of PMOS ON and PMOS OFF the mirror image of
          NMOS OFF: with $V' = V_{DD}-V$, all currents negated and $|$vth0$|$, the device equations
          take the same form. NMOS ON therefore reuses the PMOS ON pipeline and PMOS OFF reuses the
          NMOS OFF pipeline; relative errors are identical in both views.
\end{itemize}

\textbf{LightGBM hyperparameters} (all cases): Huber loss ($\alpha=0.005$), learning rate 0.08,
31 leaves, at least 20 samples per leaf, linear trees ($\lambda_\text{lin}=0.1$), 80\,\% row and
90\,\% feature subsampling, $L_2=0.01$, up to 45{,}000 trees (20{,}000 for the OFF cases) with early
stopping after 300 rounds without validation improvement, and 3 seeds (42, 7, 123) averaged.

\begin{table}[H]
\centering
\begin{tabular}{lp{11cm}}
\hline
Case & Final model \\
\hline
NMOS ON & 4 physics versions (A: best channel fit; B: + gate-to-body tunnelling; C: B + velocity
          saturation; D: B + inversion-charge gate tunnelling) $\times$ 3 seeds; per-current blend of
          the versions with weights chosen on the validation set. \\
NMOS OFF & 1 physics version $\times$ 3 seeds. \\
PMOS ON & 2 physics versions (A, B: + gate-to-body tunnelling) $\times$ 3 seeds; average of A and B,
          $I_s$ from A. \\
PMOS OFF & NMOS OFF pipeline on mirrored data, PMOS nominal values, $\times$ 3 seeds. \\
\hline
\end{tabular}
\caption{Chosen model per case.}
\end{table}

\subsection{Model development and selection}
Table~\ref{tab:ablation} compares, on the validation set, the fitted physics formulas alone with the
final physics + LightGBM models. The ML correction reduces the median error by one to two orders of
magnitude.

\begin{table}[H]
\centering
\gentable{tab_models_ablation}
\caption{Validation set: physics baseline alone vs final physics + LightGBM model.}
\label{tab:ablation}
\end{table}

For PMOS OFF two model families were trained on the same split: an ensemble of tanh MLPs
(L-BFGS, 2--4 hidden layers) that predicts $\log I_g$, $\log I_b$ and the antisymmetric part
$c=(I_d-I_s)/2$, and the physics + LightGBM model. The MLP rebuilds $I_d=-(I_g+I_b)/2+c$, which
assumes an equal split of the gate current between drain and source; near the zero crossing of
$I_d$ this becomes a difference of two nearly equal terms and the relative error explodes. The
physics + LightGBM model is better on every current on the validation set (Table~\ref{tab:pmosoff})
and was selected.

\begin{table}[H]
\centering
\gentablenat{tab_pmos_off_selection}
\caption{PMOS OFF model selection: validation mean accuracy (\%).}
\label{tab:pmosoff}
\end{table}

\subsection{Test-set results of the chosen models}
Table~\ref{tab:models} gives the test-set median and 95th-percentile error of each chosen model.
The lowest mean accuracy is \ModelMinAcc\,\% (\ModelMinAccWhere); every current of every case is above
99\,\%. Figures~\ref{fig:par_non}--\ref{fig:par_poff} show predicted vs SPICE currents.

\begin{table}[H]
\centering
\gentable{tab_models_test}
\caption{Chosen models, test set (1{,}500 samples per case).}
\label{tab:models}
\end{table}

\begin{figure}[H]
\centering
\includeimage[0.7\linewidth]{figures/parity_nmos_on.png}{NMOS ON parity}
\caption{NMOS ON: predicted vs SPICE, test set.}
\label{fig:par_non}
\end{figure}

\begin{figure}[H]
\centering
\includeimage[0.7\linewidth]{figures/parity_nmos_off.png}{NMOS OFF parity}
\caption{NMOS OFF: predicted vs SPICE, test set.}
\end{figure}

\begin{figure}[H]
\centering
\includeimage[0.7\linewidth]{figures/parity_pmos_on.png}{PMOS ON parity}
\caption{PMOS ON: predicted vs SPICE, test set.}
\end{figure}

\begin{figure}[H]
\centering
\includeimage[0.7\linewidth]{figures/parity_pmos_off.png}{PMOS OFF parity}
\caption{PMOS OFF: predicted vs SPICE, test set.}
\label{fig:par_poff}
\end{figure}

% ===========================================================================
\section{Gate-Level Results}
% ===========================================================================
\subsection{Method}
\begin{itemize}
    \item \textbf{Corners.} 200 random corners (seed 2026); a corner is one draw of $V_{DD}$,
          temperature, the NMOS process parameters and, separately, the PMOS process parameters
          (distributions of Section~III). All NMOS of a gate share one draw, all PMOS another.
    \item \textbf{Gates.} Static CMOS INV, NAND2 (internal node $x$ between the series NMOS), NOR2
          (internal node $y$ between the series PMOS), AND2 = NAND2 + INV and OR2 = NOR2 + INV,
          simulated as complete gates. Inputs at exactly 0 or $V_{DD}$; 18 cases $\times$ 200 corners =
          3{,}600 SPICE simulations, used for evaluation only.
    \item \textbf{Leakage.} Current delivered by the $V_{DD}$ source plus every input source held at
          logic 1 (static power$/V_{DD}$): $I_{leak} = -i(V_{DD}) - \sum_{in=1} i(V_{in})$.
    \item \textbf{ML prediction.} For each transistor, $V_d$ and $V_s$ come from the recorded SPICE
          node voltages (output and internal nodes), $V_g$ is the logic level of its gate node (which
          also selects ON/OFF), $V_b$ is 0 (NMOS) or $V_{DD}$ (PMOS). The four models give its four
          currents; the gate leakage is the sum of the predicted currents at the terminals on
          $V_{DD}$-level nodes. Re-simulating every transistor alone at the same voltages reproduces
          the gate's SPICE leakage exactly, so any remaining error is model error.
    \item \textbf{Rail clamp.} Inside a gate the source of an OFF transistor is exactly on a rail
          (0 for NMOS, $V_{DD}$ for PMOS), a point the training data never contains (closest
          7--31\,$\mu$V) and where the OFF-state physics term $\ln V$ is undefined. Such a terminal is
          moved 30\,$\mu$V off the rail. The value was chosen on 400 \emph{new} single-transistor SPICE
          samples per OFF case with the source exactly on the rail (not on the gate results): there the
          OFF models reach 99.34\,\% (NMOS) and 99.41\,\% (PMOS) mean accuracy for $I_d$, the current that
          carries the OFF transistor's share of the gate leakage.
\end{itemize}

\subsection{Results}
Over all \textbf{3{,}600} gate simulations the ML leakage has a mean accuracy of
\textbf{\GateMeanAcc\,\%}, median error \GateMedErr\,\%, 95th-percentile error \GatePninetyfive\,\% and
bias \GateBias\,\% (Table~\ref{tab:gate}); \GateNinetyNineCount{} of the 18 cases reach 99\,\% mean
accuracy. The most accurate case is \GateBest; the least accurate is \GateWorst{}
(mean accuracy \GateWorstAcc\,\%, p95 \GateWorstPninetyfive\,\%). By gate type, \GateBestGate{} is the
most and \GateWorstGate{} the least accurate (Table~\ref{tab:gatebygate}).

The gate-level accuracy is therefore below the 99\,\% that every transistor model reaches on its own
test set. Section~VII shows why: inside a gate the ON transistors work at an almost zero $V_{ds}$,
a corner of the input space the ON-state training data covers only sparsely.

\begin{table}[H]
\centering
\gentable{tab_gate}
\caption{Gate results: SPICE vs ML leakage for every gate and input combination (200 corners each).
Bias = median of (ML$-$SPICE)/SPICE.}
\label{tab:gate}
\end{table}

\begin{table}[H]
\centering
\gentablenat{tab_gate_by_gate}
\caption{Gate results per gate type.}
\label{tab:gatebygate}
\end{table}

\begin{figure}[H]
\centering
\includeimage[0.72\linewidth]{figures/gate_parity.png}{Gate parity plot}
\caption{ML vs SPICE leakage for all 3{,}600 gate simulations (log axes).}
\label{fig:gate_parity}
\end{figure}

\begin{figure}[H]
\centering
\includeimage[0.95\linewidth]{figures/gate_parity_by_gate.png}{Gate parity per gate}
\caption{ML vs SPICE leakage per gate type.}
\end{figure}

\begin{figure}[H]
\centering
\includeimage[0.95\linewidth]{figures/gate_error_by_case.png}{Gate error by case}
\caption{Error distribution for each gate and input combination.}
\end{figure}

% ===========================================================================
\section{Test-Circuit Results: 3-bit Carry-Lookahead Adder}
% ===========================================================================
\subsection{Circuit and method}
The adder has 7 inputs ($A_{2..0}$, $B_{2..0}$, $C_{in}$). It computes $P_i=A_i\oplus B_i$,
$G_i=A_iB_i$, the lookahead carries $C_1$--$C_3$ with AND2/OR2 gates and $S_i=P_i\oplus C_i$; each XOR
is built from four NAND2. In total 24 NAND2, 9 AND2 and 6 OR2 (39 gates, 186 transistors). The
complete netlist was simulated in SPICE for all 128 input vectors at the same 200 corners (25{,}600
simulations); the adder was checked to add correctly for every vector.

For the ML prediction the logic value at the inputs of every gate is computed for each vector, the
corresponding gate model (Part B, which calls the transistor models) gives that gate's leakage, and
the gate leakages are summed. This assumes every gate input sits at exactly 0 or $V_{DD}$ and
ignores loading between gates.

\subsection{Results}
Over all 25{,}600 simulations the hierarchical ML leakage has a mean accuracy of
\textbf{\ClaMeanAcc\,\%}, median error \ClaMedErr\,\%, p95 error \ClaPninetyfive\,\% and bias
\ClaBias\,\% (Table~\ref{tab:cla}, Fig.~\ref{fig:cla_parity}).

\begin{table}[H]
\centering
\gentable{tab_cla}
\caption{3-bit CLA: SPICE vs hierarchical ML leakage, grouped by the number of inputs at logic 1.}
\label{tab:cla}
\end{table}

\begin{figure}[H]
\centering
\includeimage[0.72\linewidth]{figures/cla_parity.png}{CLA parity plot}
\caption{3-bit CLA: hierarchical ML vs full-netlist SPICE leakage (log axes).}
\label{fig:cla_parity}
\end{figure}

\subsection{Comparison with the gate level}
The circuit error has two sources, separated in Table~\ref{tab:decomp}: the \emph{composition error}
of summing ideal gate leakages (exact SPICE gate leakages vs the full netlist; median
\ClaCompMedErr\,\%, mean accuracy \ClaCompMeanAcc\,\%) and the \emph{model error} of the ML gate
predictions (median \ClaModelMedErr\,\%). The composition error comes from the assumptions above:
internal gate inputs are not exactly at the rails and neighbouring gates load each other; on its own
it would limit the circuit to \ClaCompMeanAcc\,\% mean accuracy. Here, however, the model error
dominates. The circuit (\ClaMeanAcc\,\%) ends up slightly \emph{more} accurate than the isolated
gates (\GateMeanAcc\,\%): its leakage is the sum of 39 gate leakages, so the gate errors partly
average out, and its p95 error (\ClaPninetyfive\,\%) is lower than the gate p95
(\GatePninetyfive\,\%). The common bias of about $-1$\,\% does not average out and sets the median
error of both.

\begin{table}[H]
\centering
\gentable{tab_cla_decomp}
\caption{Decomposition of the CLA error.}
\label{tab:decomp}
\end{table}

\begin{figure}[H]
\centering
\includeimage[0.72\linewidth]{figures/cla_parity_gatesum.png}{CLA SPICE gate-sum parity}
\caption{Composition error alone: sum of SPICE gate leakages vs full-netlist SPICE.}
\end{figure}

% ===========================================================================
\section{Error Analysis}
% ===========================================================================
\subsection{Transistor level}
\begin{itemize}
    \item \textbf{Zero crossing of $I_d$ / $I_s$.} The largest relative errors occur where the drain
          or source current changes sign: in the OFF state the subthreshold channel current and the
          gate-edge tunnelling current flow in opposite directions at one terminal, so the terminal
          current can be orders of magnitude smaller than either part. Any small absolute error is a
          large relative error there. This dominates the maximum errors of NMOS OFF and PMOS OFF and
          is why the ratio target $I/(C+F)$ was used.
    \item \textbf{Body current in the OFF state.} $I_b$ sits close to $10^{-14}$\,A (junction floor) in
          most OFF samples, and its relative variation is small, so its linear $R^2$ is low even
          though its mean accuracy is above 99\,\%. It is negligible for gate leakage.
    \item \textbf{ON-state gate and body currents.} NMOS ON and PMOS ON have the highest
          95th-percentile errors (1.4--3.5\,\%) in $I_g$ and $I_b$, where gate-to-channel and
          gate-to-body tunnelling change with the inversion charge along the channel.
\end{itemize}

\subsection{Gate and circuit level}
\textbf{The ON transistors cause most of the gate error.} Table~\ref{tab:tingate} compares every
transistor's predicted terminal currents inside the gates with a SPICE simulation of that transistor
alone at the same voltages. The ON transistors account for \OnShare\,\% of the absolute gate-leakage
error (PMOS ON \PmosOnShare\,\%, NMOS ON \NmosOnShare\,\%), although their models reach 99.1--99.6\,\%
on their own test sets.

The reason is a mismatch of operating points. In a static gate an ON transistor connects the output
(or an internal node) to a rail and carries only the leakage of the OFF network, so its $|V_{ds}|$
is tiny: median \OnVdsMedian\,mV, 90\,\% below \OnVdsNinety\,mV. The ON-state datasets draw $V_d$ and
$V_s$ uniformly in $[0, V_{DD}]$, so only \OnTrainSmallVds\,\% of their samples have $|V_{ds}|<10$\,mV
and almost none are in the sub-millivolt range. There the channel current is nearly zero and the
terminal currents are dominated by gate tunnelling, which the models have rarely seen. The OFF
transistors, whose drain-source voltage in a gate is large and well covered by training, predict
their $I_d$ to a median of 0.8--0.9\,\%.

\begin{table}[H]
\centering
\gentable{tab_transistor_in_gate}
\caption{Transistor currents inside the gates, ML vs single-transistor SPICE at the same voltages
(median / p95 error in \%), and each case's share of the absolute gate-leakage error.}
\label{tab:tingate}
\end{table}

The worst gate case is \GateWorst{}: both NMOS are OFF in series and both PMOS are ON in parallel
with almost zero $V_{ds}$, so two ON-state errors add up. Temperature has little influence: the median
gate error is \GateErrHot\,\% for $T \ge 85\,^\circ$C and \GateErrCold\,\% for $T\le0\,^\circ$C (CLA:
\ClaErrHot\,\% and \ClaErrCold\,\%).

\begin{figure}[H]
\centering
\includeimage[0.95\linewidth]{figures/gate_error_vs_corner.png}{Gate error vs corner}
\caption{Gate error vs corner conditions.}
\end{figure}

\begin{figure}[H]
\centering
\includeimage[0.95\linewidth]{figures/cla_error_vs_corner.png}{CLA error vs corner}
\caption{CLA error vs corner conditions.}
\end{figure}

\begin{figure}[H]
\centering
\includeimage[0.72\linewidth]{figures/cla_error_by_ones.png}{CLA error by input pattern}
\caption{CLA error vs the number of inputs at logic 1.}
\end{figure}

% ===========================================================================
\section{Limitations and Improvements}
% ===========================================================================
\begin{itemize}
    \item \textbf{Training distribution vs gate operating points (main limitation).} The ON models
          are trained on uniformly drawn $V_d$, $V_s$, but in gates they operate at
          $|V_{ds}|\approx0$ (median \OnVdsMedian\,mV). The most effective improvement is to add
          ON-state training samples with $|V_{ds}|$ from $10^{-6}$ to $10^{-2}$\,V (for example
          log-uniform), and OFF-state samples with a terminal exactly on the rail, and retrain. The
          project document allows the dataset size to be increased for this purpose.
    \item \textbf{Training data size.} Each model sees 7{,}000 samples. The 95th-percentile errors of
          the ON-state models (1.4--3.5\,\%) would also fall with more samples.
    \item \textbf{Zero crossings.} Relative error is ill-defined where a terminal current passes
          through zero. Sampling more points near $V_d\approx V_s$ and near the drain-current zero
          crossing of the OFF state, or reporting a floor-normalised error, would address this.
    \item \textbf{Composition assumptions.} The circuit model assumes rail-to-rail gate inputs and no
          loading. Passing the predicted (not ideal) output voltage of each gate to the next, or
          adding the internal-node voltages from a fast DC solve, should reduce the
          composition error.
    \item \textbf{Run time.} Each physics fit takes a few minutes and each LightGBM model
          4--9 minutes on one core. Predicting all 16{,}800 gate transistors takes about five minutes
          on 12 threads, most of it for the 12 large NMOS ON models.
\end{itemize}

% ===========================================================================
\section{Reproducibility}
% ===========================================================================
Every result is regenerated from the model file by scripts with fixed seeds. \texttt{run\_all.sh}
runs the stages in order:

\begin{table}[H]
\centering
\begin{tabular}{lp{10.5cm}}
\hline
Stage & Scripts \\
\hline
data & \texttt{dataset/<CASE>/01\_*\_inputs.py}, \texttt{02\_*\_simulate.py} (seed 42) \\
datacheck & \texttt{dataset/03\_data\_check.py}, \texttt{04\_example\_netlists.py} \\
train & \texttt{ml\_model/train\_parallel.py} (NMOS ON, PMOS ON), \texttt{train\_final.py} (NMOS OFF, PMOS OFF) \\
check & \texttt{ml\_model/summarize\_models.py}, \texttt{gate\_circuit/00\_check\_models.py} \\
spice & \texttt{gate\_circuit/01}--\texttt{04} (corners seed 2026, gate and CLA SPICE) \\
gates & \texttt{gate\_circuit/05}--\texttt{08} (ML gate and CLA leakage, tables, plots) \\
collect & \texttt{collect\_deliverables.py} \\
report & \texttt{report/make\_report\_assets.py} \\
\hline
\end{tabular}
\caption{Pipeline stages.}
\end{table}

\end{document}
